\section*{Capítulo \ref{ch:b3:palaeomagnetism} --- Geomagnetismo e paleomagnetismo}

\begin{solution}{exo:b3:palaeomagnetism:1}
$\tan I = 2\tan\ang{30} = 1.1547$, logo $I = \ang{49.1}$. Inversamente, $\tan\lambda
= \tfrac12\tan(\ang{-20}) = -0.1820$, logo $\lambda = \ang{-10.3}$, ou seja
\ang{10.3}\,S.
\end{solution}

\begin{solution}{exo:b3:palaeomagnetism:2}
$H = \sqrt{20000^2 + 2000^2} = \qty{20100}{nT}$;
$F = \sqrt{20000^2 + 2000^2 + 40000^2} = \qty{44770}{nT}$. A
declinação magnética é $\arctan(-2000/20000) = \ang{-5.7}$, \ang{5.7}~O do norte;
a inclinação $\arctan(40000/20100) = \ang{63.3}$, para baixo.
\end{solution}

\begin{solution}{exo:b3:palaeomagnetism:3}
$B_0 = -g_1^0 = \qty{29350}{nT}$ no equador, e $2B_0 = \qty{58700}{nT}$ nos
polos. $m = 4\pi a^3B_0/\mu_0 = a^3B_0/\num{e-7} = \num{2.5862e20}\times
\num{2.935e-5}/\num{e-7} = \qty{7.59e22}{A.m^2}$, um pouco menor que o
\qty{7.69e22}{A.m^2} do dipolo completo, que acrescenta $g_1^1$ e $h_1^1$.
\end{solution}

\begin{solution}{exo:b3:palaeomagnetism:4}
(a) Uma \omterm{def:b3:palaeomagnetism:nrm}{magnetização termorremanente}, que registra o campo à medida que a lava
resfriava ao passar pelas \omterm{def:b3:palaeomagnetism:blocking}{temperaturas de bloqueio} de seus grãos. (b) Uma
\omterm{def:b3:palaeomagnetism:nrm}{magnetização remanente detrítica}, que registra o campo durante o tempo que o silte levou para assentar e
ser fixado pela compactação. (c) Uma \omterm{def:b3:palaeomagnetism:nrm}{magnetização remanente química}, que registra o
campo durante o crescimento da hematita, no soterramento inicial e, portanto, um pouco
depois da deposição. (d) Uma remanência viscosa na direção do campo atual, somada à
TRM original pelos grãos de menor tempo de relaxação; ela é removida pelas
primeiras etapas da desmagnetização.
\end{solution}

\begin{solution}{exo:b3:palaeomagnetism:5}
Os vetores unitários $(\cos I\cos D,\ \cos I\sin D,\ \sin I)$ são
\begin{gather*}
  (0.6545, 0.1391, 0.7431),\quad (0.6670, 0.2695, 0.6947),\\
  (0.5235, 0.1501, 0.8387),\quad (0.5985, 0.0629, 0.7986).
\end{gather*}
A soma é $(2.4435, 0.6216, 3.0751)$, de comprimento $R = 3.9766$. Média:
$D = \arctan(0.6216/2.4435) = \ang{14.3}$, $I = \arcsin(3.0751/3.9766) =
\ang{50.7}$. $k = 3/(4 - 3.9766) = 128$. Com $20^{1/3} - 1 = 1.7144$,
$\cos\alpha_{95} = 1 - (0.0234/3.9766)\times1.7144 = 0.98991$, logo
$\alpha_{95} = \ang{8.1}$: com quatro espécimes o cone permanece amplo, por mais
concentrado que seja o grupo.
\end{solution}

\begin{solution}{exo:b3:palaeomagnetism:6}
$\lambda = \arctan(\tfrac12\tan\ang{50}) = \ang{30.79}$, e $p = \ang{59.21}$.
$\sin\lambda_p = \sin\ang{40}\cos\ang{59.21} + \cos\ang{40}\sin\ang{59.21}
\cos\ang{20} = 0.9474$, $\lambda_p = \ang{71.34}$\,N. $\sin\beta =
\sin\ang{59.21}\sin\ang{20}/\cos\ang{71.34} = 0.9184$, $\beta = \ang{66.66}$.
Como $\cos p = 0.5119 < \sin\ang{40}\times0.9474 = 0.6090$, o segundo caso
aplica-se: $\varphi_p = 10 + 180 - 66.66 = \ang{123.34}$\,L. O polo está em
\ang{71.3}\,N, \ang{123.3}\,L.
\end{solution}

\begin{solution}{exo:b3:palaeomagnetism:7}
(a) $(7.69 - 8.32)/8.32 = \qty{-7.6}{\percent}$ em 125 anos, ou
\qty{-6.1}{\percent} por século. (b) $29350/12.6 = 2329$, cerca de 2300 anos.
Não se trata de previsão: a própria taxa muda de um modelo de cinco anos para o
seguinte, o momento já esteve tanto acima quanto abaixo do atual no passado sem
inverter-se, e o campo se inverte em intervalos irregulares que nenhuma taxa
prevê.
\end{solution}

\begin{solution}{exo:b3:palaeomagnetism:8}
Entre 20 e \qty{300}{\celsius} o vetor perdeu $(6.4, 0.0, 13.6)$: uma
componente de \qty{15.0}{mA/m}, com $D = \ang{0.0}$ e
$I = \arctan(13.6/6.4) = \ang{64.8}$. Acima de \qty{300}{\celsius}, o vetor perdeu $(25.8, 18.1, 24.6)$:
\qty{40.0}{mA/m}, $D = \arctan(18.1/25.8) = \ang{35.1}$, $I =
\arctan(24.6/31.5) = \ang{38.0}$. A segunda, que vai à origem, é a
ChRM; a primeira, ao longo de uma direção íngreme voltada para o norte, removida a
baixa temperatura, é uma sobreimpressão viscosa ao longo do campo recente.
\end{solution}

\begin{solution}{exo:b3:palaeomagnetism:9}
$\lambda = \arctan(\tfrac12\tan(\ang{-35})) = \ang{-19.3}$. Para $I = \ang{-40}$
e \ang{-30}: \ang{-22.8} e \ang{-16.1}, um intervalo de \ang{6.7}.
$p = \ang{90} - \lambda = \ang{109.3}$, $\cos^2p = 0.1096$:
$\mathrm dp = 2.5\times(1 + 0.329) = \ang{3.3}$ e $\mathrm dm =
5\times\sin\ang{109.3}/\cos\ang{35} = \ang{5.8}$.
\end{solution}

\begin{solution}{exo:b3:palaeomagnetism:10}
(a) $\tan\lambda = \tfrac12$, $\lambda = \ang{26.6}$. $F$ vale $B_0$ no
equador e $2B_0$ nos polos: uma razão de 2. (b) $\tan I = 2\tan\lambda$
fornece $(1 + \tan^2I)\,\mathrm dI = 2(1 + \tan^2\lambda)\,\mathrm d\lambda$, logo
\[
  \frac{\mathrm dI}{\mathrm d\lambda} = \frac{2/\cos^2\lambda}{1 + 4\tan^2\lambda}
  = \frac{2}{\cos^2\lambda + 4\sin^2\lambda} = \frac{2}{1 + 3\sin^2\lambda}.
\]
Um erro $\delta I$ fornece $\delta\lambda = \tfrac12(1 + 3\sin^2\lambda)\,\delta I$,
de $\tfrac12\delta I$ no equador a $2\,\delta I$ nos polos: o
maior erro está perto dos polos. (c) Perto do equador, a inclinação varia
duas vezes mais rápido que a latitude, logo um dado erro de inclinação custa metade
da latitude; perto dos polos, a inclinação quase não varia com a latitude.
\end{solution}

\begin{solution}{exo:b3:palaeomagnetism:11}
(a) $\cos p = \sin(\ang{-35})\sin\ang{70} + \cos\ang{35}\cos\ang{70}\cos\ang{150}
= -0.5390 - 0.2426 = -0.7816$, $p = \ang{141.41}$. Então $\tan I_x = 2\cot p =
-2.506$, $I_x = \ang{-68.2}$, e $\lambda_x = \ang{90} - p = \ang{-51.4}$.
$\cos D_x = (\sin\ang{70} - \sin(\ang{-35})\cos p)/(\cos\ang{35}\sin p) =
0.9617$, e, como o polo fica a leste do meridiano do sítio ($\sin\ang{150} >
0$), $D_x = \ang{15.9}$. (b) $\lambda_o = \arctan(\tfrac12\tan(\ang{-60})) =
\ang{-40.9}$. $R = 30.0 - 15.9 = \ang{14.1}$: uma rotação horária do bloco
em torno de um eixo vertical. $P = \lambda_x - \lambda_o = -51.4 + 40.9 =
\ang{-10.5}$: a formação ficava \ang{10.5} mais ao norte do que a placa estável
indica que deveria; logo o bloco se deslocou desde então cerca de \ang{10.5}, algo como
\qty{1170}{km} para o sul em relação à placa.
\end{solution}

\begin{solution}{exo:b3:palaeomagnetism:12}
(a) Com $\mathrm d\Omega = 2\pi\sin\theta'\,\mathrm d\theta'$,
\[
  \int_0^\theta\frac{\kappa e^{\kappa\cos\theta'}}{4\pi\sinh\kappa}\,2\pi\sin\theta'\,
  \mathrm d\theta' = \frac{e^\kappa - e^{\kappa\cos\theta}}{2\sinh\kappa}
  = \frac{e^\kappa - e^{\kappa\cos\theta}}{e^\kappa - e^{-\kappa}} .
\]
(b) Para $\kappa$ grande, $e^{-\kappa}$ é desprezível; dividindo por $e^\kappa$
obtém-se $1 - e^{-\kappa(1 - \cos\theta)} \approx 1 - e^{-\kappa\theta^2/2}$. Isso
vale $1 - e^{-1} = \qty{63}{\percent}$ quando $\kappa\theta^2/2 = 1$, $\theta =
\sqrt{2/\kappa}$ radianos, ou $\ang{81}/\sqrt\kappa$. (c) $81/\sqrt{50} =
\ang{11.5}$; a fórmula exata dá \qty{63.1}{\percent}.
\end{solution}

\begin{solution}{pb:b3:palaeomagnetism:1}
\textbf{1.} A hematita em uma camada vermelha pode crescer a partir de óxidos e hidróxidos de ferro
muito depois da deposição do sedimento, gerando uma remanência química dessa
idade posterior. Um teste de dobramento positivo, caso as camadas tenham sido dobradas logo após a
deposição, ou um teste de inversão positivo em camadas próximas entre si com ambas as
polaridades, mostraria que a remanência é antiga.

\textbf{2.} $(\cos\ang{37}\cos\ang{50}, \cos\ang{37}\sin\ang{50}, \sin\ang{37})
= (0.5134, 0.6118, 0.6018)$.

\textbf{3.} $R = \sqrt{3.8943^2 + 2.6500^2 + 3.5717^2} = 5.9114$. $D =
\arctan(2.6500/3.8943) = \ang{34.2}$, $I = \arcsin(3.5717/5.9114) = \ang{37.2}$.

\textbf{4.} $k = 5/(6 - 5.9114) = 56.5$. Com $20^{1/5} - 1 = 0.8206$,
$\cos\alpha_{95} = 1 - (0.0886/5.9114)\times0.8206 = 0.98771$, $\alpha_{95} =
\ang{9.0}$.

\textbf{5.} B: $k = 6/(7 - 6.8979) = 58.8$; $20^{1/6} - 1 = 0.6475$,
$\cos\alpha_{95} = 1 - (0.1021/6.8979)\times0.6475 = 0.99042$, $\alpha_{95} =
\ang{7.9}$. C: $k = 7/(8 - 7.8707) = 54.1$; $20^{1/7} - 1 = 0.5340$,
$\cos\alpha_{95} = 1 - (0.1293/7.8707)\times0.5340 = 0.99123$, $\alpha_{95} =
\ang{7.6}$.

\textbf{6.} B tem o maior $k$, o agrupamento mais apertado; C tem o menor
$\alpha_{95}$, a média mais bem definida, por ter o maior número de sítios. Os três
são próximos: as três formações estão bem agrupadas.

\textbf{7.} $\lambda_A = \arctan(\tfrac12\tan\ang{37.17}) = \ang{20.76}$\,N,
$\lambda_B = \ang{32.15}$\,N, $\lambda_C = \ang{45.33}$\,N.

\textbf{8.} $p = 90 - 20.76 = \ang{69.24}$. $\sin\lambda_p =
\sin\ang{47}\cos\ang{69.24} + \cos\ang{47}\sin\ang{69.24}\cos\ang{34.23} =
0.2593 + 0.5272 = 0.7865$: $\lambda_p = \ang{51.86}$\,N.

\textbf{9.} $\sin\beta = \sin\ang{69.24}\sin\ang{34.23}/\cos\ang{51.86} =
0.8516$, $\beta = \ang{58.40}$. Como $\cos p = 0.3545 < \sin\ang{47}\times
0.7865 = 0.5752$, aplica-se o segundo caso, $\cos\beta < 0$:
$\varphi_p = 15.00 + 180 - 58.40 = \ang{136.60}$\,L.

\textbf{10.} $\cos^2p = 0.1257$: $\mathrm dp = 4.50\times1.377 = \ang{6.2}$;
$\mathrm dm = 8.99\times\sin\ang{69.24}/\cos\ang{37.17} = \ang{10.5}$.

\textbf{11.} De \ang{51.9}\,N em \qty{60}{Ma} até \ang{66.5}\,N em
\qty{30}{Ma} e \ang{87.1}\,N em \qty{4}{Ma}, sobre quase o mesmo meridiano
(\ang{136}--\ang{139}\,L): um caminho em direção ao polo geográfico, alcançado dentro
de seu cone pela formação mais jovem.

\textbf{12.} A camada de cinza, a \qty{22}{m}, situa-se na segunda zona normal,
19.7--\qty{24.6}{m}; com idade próxima de \qty{3.2}{Ma}, essa zona é C2An.2n
(3.207--3.116). Então a zona 4.2--\qty{14.8}{m} é C2An.3n, a zona
29.2--\qty{54.0}{m} é C2An.1n, a base invertida pertence ao cron abaixo de
C2An, e o topo invertido ao cron acima dele.

\textbf{13.} $(19.7 - 4.2)/(3.596 - 3.207) = 15.5/0.389 = \qty{39.8}{m/Myr}$;
$(54.0 - 24.6)/(3.116 - 2.581) = 29.4/0.535 = \qty{55.0}{m/Myr}$.

\textbf{14.} Base: $3.596 + 4.2/39.8 = \qty{3.70}{Ma}$. Topo: $2.581 -
10.0/55.0 = \qty{2.40}{Ma}$.

\textbf{15.} Somente os limites de inversão são datados, a poucas centenas de milhares de anos
de distância entre si; entre eles, assume-se que a taxa é constante. Uma mudança na subsidência, no
suprimento de sedimentos ou no clima, ou uma lacuna na deposição, a alteraria, e um
hiato poderia até eliminar uma zona curta e enganar o ajuste.

\textbf{16.} $\cos p = \sin\ang{47}\sin\ang{78} + \cos\ang{47}\cos\ang{78}
\cos\ang{175} = 0.7154 - 0.1413 = 0.5741$, $p = \ang{54.96}$. $\tan I_x =
2\cot p = 1.4023$, $I_x = \ang{54.5}$; $\lambda_x = \ang{35.04}$\,N.

\textbf{17.} $\cos D_x = (\sin\ang{78} - \sin\ang{47}\times0.5741)/
(\cos\ang{47}\sin\ang{54.96}) = 0.99976$; o polo fica a leste do meridiano do
sítio ($\sin\ang{175} > 0$), logo $D_x = \ang{1.3}$.

\textbf{18.} \omterm{def:b3:palaeomagnetism:vgp}{Paleolatitudes} esperadas a partir de $I_x$ = \ang{54.5}, \ang{60.0} e
\ang{63.6}: \ang{35.03}, \ang{40.89} e \ang{45.21}. $P_A = 35.03 - 20.76 =
\ang{14.3}$, $P_B = 40.89 - 32.15 = \ang{8.7}$, $P_C = 45.21 - 45.33 =
\ang{-0.1}$.

\textbf{19.} $R_A = 34.2 - 1.3 = \ang{32.9}$, $R_B = 23.9 - 2.9 = \ang{21.0}$,
$R_C = 3.4 - 2.5 = \ang{0.9}$, todas no sentido horário.

\textbf{20.} $\Delta D_o = \arcsin(\sin\ang{8.99}/\cos\ang{37.17}) = \ang{11.3}$;
$\Delta D_x = \arcsin(\sin\ang{3.0}/\cos\ang{54.5}) = \ang{5.2}$; $\Delta R_A =
\sqrt{11.3^2 + 5.2^2} = \ang{12.4}$.

\textbf{21.} $R_A = \ang{32.9} \pm \ang{12.4}$ e $R_B = \ang{21.0} \pm
\ang{13.8}$ são significativas; $R_C = \ang{0.9} \pm \ang{17.9}$ não é.

\textbf{22.} Sessenta milhões de anos atrás, Kesta situava-se cerca de \ang{14} de latitude
ao sul de onde o continente estável a coloca, e fora girada cerca de \ang{33}
no sentido anti-horário em relação à sua orientação atual. Em \qty{30}{Ma} ela havia avançado
\ang{6} desses \ang{14} para o norte e girado \ang{12} desses \ang{33}
no sentido horário; em \qty{4}{Ma} alcançara seu lugar atual no continente,
sem rotação restante. Ela derivou para o norte ao longo da margem do continente, girando
no sentido horário à medida que avançava, e estava no lugar por volta de \qty{4}{Ma}.

\textbf{23.} $47.00 - 20.76 = \ang{26.24}$, ou $26.24\times111.19 =
\qty{2918}{km}$. Isso inclui o movimento do próprio continente estável;
em relação a ele, Kesta se deslocou $P_A = \ang{14.3}$.

\textbf{24.} $I = \ang{28.18}$ e \ang{46.16} dão $\lambda = \ang{15.00}$ e
\ang{27.50}: uma semilargura de \ang{6.25}, e um desvio-padrão de
$6.25/1.96 = \ang{3.19}$, ou \qty{355}{km}. Desde a formação A, Kesta derivou
para o norte em
\textbf{\qty[text-series-to-math=true]{2918 \pm 355}{km}}.
\end{solution}
